\documentclass[../../main2.tex]{subfiles}

\begin{document}

\chapter{从贡献表到未来走向}
\label{ch:forecast}

\begin{motivation}
	上一章遗留的问题：四条通道定义完毕，分账公式在手——但那是对\textbf{过去}的分解。手里攥着四个通道的贡献数，怎么用它说出「接下来会怎样」？\\
	\textbf{核心动机}：把贡献表拼回时间线，完整演示一次「怎么用这本书」——这是全书唯一一次手把手的推演。然后自曝其短：失败模式清单。
\end{motivation}

\begin{logicmap}
\begin{tikzpicture}[node distance=0.85cm, every node/.style={draw, rectangle, rounded corners, align=center, font=\small}]
\node (q) {贡献表怎么变「接下来」？};
\node (s) [below=of q] {贡献表 = 状态压缩};
\node (a) [below=of s] {链条重组：从分解到预测};
\node (e) [below=of a] {完整演示：防沉迷新规};
\node (f) [below=of e] {失败模式清单};
\draw[-{Stealth}] (q) -- (s);
\draw[-{Stealth}] (s) -- (a);
\draw[-{Stealth}] (a) -- (e);
\draw[-{Stealth}] (e) -- (f);
\end{tikzpicture}
\end{logicmap}

\section{贡献表是一份状态压缩}

\begin{readingnote}
	你有没有想过，为什么天气预报不说「每一滴雨落在哪里」？\\
	因为预测必须\textbf{丢掉}细节，才能换来可用的答案。
\end{readingnote}

四个通道的贡献数（$C_{\text{观念}}, C_{\text{经济}}, C_{\text{政治}}, C_{\text{法律}}$）是一张表，也是一份压缩：把亿万个个体、千万条动机，压成几个能上桌的数。压缩必有丢失（第 2 章的地图类比）——丢掉个体差异、丢掉通道内部的异质性。

丢多少合适？看预测目标的尺度。预测「明年全国房价走势」，丢掉户型细节没事；预测「你家房子卖几万」，这份压缩就太粗了。\textbf{压缩的粒度要和问题的粒度配对}——这是第一个使用纪律。

\paragraph*{逻辑衔接}
	表有了，粒度对齐了。怎么把它变回「未来」？下一节把第 6 章的链条重新组装。

\section{链条重组：从分解到预测}

第 6 章说过，预测是沿因果链的概率接力：每个环节估一个「多大概率过」，乘起来是整链。现在把链上的环节从「抽象因素」换成「通道」：

\begin{enumerate}
	\item \textbf{定对象}：预测什么、多长时间、多大范围（例：一年内、某市、青少年日均使用时长）。
	\item \textbf{列通道}：把四个齿轮各写一条「当前状态 + 可能变化」（观念风向、平台营收压力、监管姿态、条文空子）。
	\item \textbf{估贡献}：对每个通道做一次无你检验——「拿掉它，结果会变多少」——用第 9 章的规矩，份额量级成对记。
	\item \textbf{排顺序}：哪个通道是短板（第 6 章：整链被最弱环节压住），哪个是杠杆点。
	\item \textbf{出断言}：「在 X 不变的前提下，Y 有较大把握发生」——前提写明，断言才可检验。
\end{enumerate}

\paragraph*{逻辑衔接}
	五步有了。空讲没用——下一节，一个真实政策，走完整的一遍。

\section{完整演示：防沉迷新规会怎么走}

选题：某平台上线「未成年人防沉迷限制」（限时、实名、宵禁）。预测对象：半年后，未成年人日均使用时长。下面的数字是\lowfreq{演示用的示意值}——不是真实预测，是展示\textbf{怎么算}；真实做预测时，每一步的数都要换上你的证据（问卷、财报、历史案例）。

\begin{mathbox}{演示推演：四通道分账（示意数）}
	\textbf{第 1 步，基准}：现状日均 3.2 小时。目标是估计「新规后降到多少」。\\
	\textbf{第 2 步，通道贡献估计}（各自做无你检验，历史案例 + 财报 + 舆情为证据）：
	\begin{itemize}
		\item 法律（$C_L$）：限时条文直接砍时段。历史同类（游戏防沉迷）降幅约 25\%——$C_L \approx -0.8$ 小时；
		\item 经济（$C_E$）：平台营收靠时长，KPI 会逼出「防沉迷松紧带」。阻力约 $+0.3$ 小时；
		\item 观念（$C_M$）：家长舆论支持、青少年「上有政策下有对策」。净约 $-0.2$ 小时；
		\item 政治（$C_P$）：监管约谈威慑，平台不敢公然放水。约 $-0.3$ 小时。
	\end{itemize}
	\textbf{第 3 步，交互项}：法律与经济互相拉扯（规则 + 稀缺 = 套利，第 14 章），交互约 $+0.2$ 小时——单挂给谁都错，按第 9 章规矩\textbf{单列}。\\
	\textbf{第 4 步，合成}：$3.2 + (-0.8) + 0.3 + (-0.2) + (-0.3) + 0.2 = 2.4$ 小时。\\
	\textbf{第 5 步，断言}：「若监管约谈力度不降（前提），半年后日均约 2.2–2.6 小时（概率断言）；短板是经济阻力与交互套利。」
\end{mathbox}

对照这个演示的三个看点：

\begin{itemize}
	\item \textbf{前提必须写出来}。「监管力度不降」一变，结论就变——可检验性藏在前提里。
	\item \textbf{交互项单列}。$+0.2$ 的套利空间是「上有政策下有对策」的定量形态，藏进任何一方都是粉饰。
	\item \textbf{短板即抓手}。想再压时长，别在 $C_L$ 上加码（已近饱和），去堵 $C_E$ 的松紧带和交互的套利口。
\end{itemize}

\begin{questionbox}
	\textit{现在你可能想反驳：「你这数字都是编的，演示有什么意义？」}\\
	\textbf{数字是示意，骨架是真的。}每个数背后都有可查的同类案例（游戏防沉迷的降幅、平台财报的时长依赖）。「禁止编造」的正确姿势不是不演示，是把演示的数\textbf{标成示意}，并告诉你真实数从哪挖。骨架你随时能套到自己的问题上。
\end{questionbox}

\paragraph*{逻辑衔接}
	推演走完一遍了。但一本诚实的书必须自曝其短——这套方法在哪里会摔？下一节清单列满。

\section{失败模式清单：方法什么时候失灵}

六条，按我判断的杀伤力排序：

\begin{enumerate}
	\item \textbf{混杂残留}：第 7 章的镊子只能逼近，不能穷尽。漏掉的混杂会让 $C$ 系统性偏高或偏低。例：把「平台改版」的功劳记给「新规」。
	\item \textbf{外推失效}（分布漂移）：从历史估的 $C$，环境一变就作废。2008 的模型、恐龙时代的生存法则，都是外推死的。
	\item \textbf{反身性}：预测公开，改变动机，预测自己作废（或自己实现）。央行预测「无风险」反而诱发冒险——下一章的主角。
	\item \textbf{尺度错配}：拿全国 $C$ 预测你家的事（上一节的粒度纪律），压缩的丢失在细尺度上会咬人。
	\item \textbf{测量误差}：「日均时长」本身就测不准——自报、埋点、口径之争。垃圾进，垃圾出。
	\item \textbf{模型误设}：漏了通道、漏了交互项、把负交互记成正。王安石的账单就是漏了「渠道」这一维。
\end{enumerate}

\begin{warningbox}
	少于三条失败模式的预测方法论，都是没认真想过的。列出来不是示弱——\textbf{每条失败模式都标注了「哪里、以什么姿势翻车」，你就知道该在哪里加倍小心}。这是本书可信度的一部分。
\end{warningbox}

\begin{reader_anticipation}
	\item 示意数字既然不是预测，真实预测的数字从哪挖？→ 财报、统计年鉴、历史案例——每类证据的用法对应附录 A.3 的「估 $C$」纪律。
	\item 预测做出来了，公开它安全吗？→ 不一定——预测会回流成新动机。第 16 章：循环的闭合。
\end{reader_anticipation}

\paragraph*{逻辑衔接}
	方法给全了，边界也亮了。还剩最后一问，也是全书的闭合点：你把预测做出来了，\textbf{预测本身}会改变世界吗？第 16 章：循环的闭合。

\begin{thinkbox}
\begin{enumerate}
	\item 用本章五步，给一个你真关心的问题出一份「概率断言」（跳槽前景、行业走势、考研上岸率）。按规矩：前提写明、交互项单列、份额量级成对。不求准，求\textbf{可检验}。
	\item \textbf{反事实}：如果防沉迷新规「只立法、不约谈」——监管威慑那一项归零——演示里的合成结果会变成多少小时？短板会挪到哪个通道？
	\item \textbf{反事实}：如果六条失败模式里只能消灭一条，你消灭哪条？为什么是它？——用「哪条的 $C$ 最大」来答，别凭感觉。
\end{enumerate}
\end{thinkbox}

\end{document}
